% GENERATED FILE. Edit canonical lessons, then run npm run generate:books.
% canonical-source-hash: fnv1a64:7d97d4cffc042281
% canonical-lessons: IT-C519-palcoscenico, IT-C519-sipario, IT-C519-platea, IT-C519-loggione, IT-C519-applauso

\chapter{A teatro}
\label{ch:it-a-teatro}
\hlchaptermodality{519}

\begin{chapteropening}
\textbf{By the end of this chapter:} \emph{I can describe a theatre and its audience.}

Complete the last lesson of chapter 519: I can describe a theatre and its audience.
\end{chapteropening}

\section[il palcoscenico]{il palcoscenico — stage}

\label{lesson:IT-C519-palcoscenico}



\begin{quote}
Before the new one: say the Italian for verse, then the Italian for to hum.
\end{quote}

\subsection*{You'll want to know: il palcoscenico}
\textbf{il palcoscenico} — \textquotedblleft{}stage\textquotedblright{}.

\begin{quote}
Gli attori salgono sul palcoscenico.
\end{quote}

\begin{quote}
\emph{The actors go up on stage.}
\end{quote}

Often shortened to \emph{il palco}.

\begin{cousinweb}
\emph{palcoscenico} is built from \emph{palco} \textquotedblleft{}platform\textquotedblright{}, usually traced to a Germanic word for a beam (EtymDB lists it beside German \emph{Balken}), and \emph{scenico} \textquotedblleft{}of the stage\textquotedblright{}, from Latin \emph{scaena}, from Greek \emph{skēn\'{\={e}}} \textquotedblleft{}a tent, a stage\textquotedblright{} (EtymDB).
\end{cousinweb}

\begin{grammarlens}[title={What you've built}]
One more word for telling, explaining and describing at B1.
\end{grammarlens}

\subsection*{Your turn}
\begin{itemize}
\raggedright
  \item \emph{Say it:} \emph{il palcoscenico}
  \item \emph{Say it:} \emph{il palcoscenico}, once more
  \item \emph{Say it:} a sentence of your own with \emph{il palcoscenico}
  \item \emph{Recall:} say the Italian for verse, then the Italian for to hum, then say \emph{il palcoscenico} again
\end{itemize}

\begin{tcolorbox}[breakable,colback=teal!4,colframe=teal!35!black,title={Before you move on}]
What does il palcoscenico mean? (\textquotedblleft{}Stage\textquotedblright{}.) Say it once more.
\end{tcolorbox}
\section[il sipario]{il sipario — curtain (in a theatre)}

\label{lesson:IT-C519-sipario}



\begin{quote}
Before the new one: say the Italian for to hum, then the Italian for stage.
\end{quote}

\subsection*{You'll want to know: il sipario}
\textbf{il sipario} — \textquotedblleft{}curtain (in a theatre)\textquotedblright{}.

\begin{quote}
Si alza il sipario e comincia lo spettacolo.
\end{quote}

\begin{quote}
\emph{The curtain rises and the show begins.}
\end{quote}

Only for the theatre; at home, \emph{la tenda}.

\begin{cousinweb}
\emph{sipario} comes from Latin \emph{sīparium} \textquotedblleft{}a stage curtain\textquotedblright{} (EtymDB), from Greek \emph{síparos} \textquotedblleft{}a topsail\textquotedblright{}.
\end{cousinweb}

\begin{grammarlens}[title={What you've built}]
One more word for telling, explaining and describing at B1.
\end{grammarlens}

\subsection*{Your turn}
\begin{itemize}
\raggedright
  \item \emph{Say it:} \emph{il sipario}
  \item \emph{Say it:} \emph{il sipario}, once more
  \item \emph{Say it:} a sentence of your own with \emph{il sipario}
  \item \emph{Recall:} say the Italian for to hum, then the Italian for stage, then say \emph{il sipario} again
\end{itemize}

\begin{tcolorbox}[breakable,colback=teal!4,colframe=teal!35!black,title={Before you move on}]
What does il sipario mean? (\textquotedblleft{}Curtain (in a theatre)\textquotedblright{}.) Say it once more.
\end{tcolorbox}
\section[la platea]{la platea — stalls; audience}

\label{lesson:IT-C519-platea}



\begin{quote}
Before the new one: say the Italian for stage, then the Italian for curtain.
\end{quote}

\subsection*{You'll want to know: la platea}
\textbf{la platea} — \textquotedblleft{}stalls; audience\textquotedblright{}.

\begin{quote}
Abbiamo i posti in platea, in quinta fila.
\end{quote}

\begin{quote}
\emph{We have seats in the stalls, in the fifth row.}
\end{quote}

Also the audience as a whole.

\begin{cousinweb}
\emph{platea} is a learned borrowing from Latin \emph{platēa} \textquotedblleft{}a broad street, a courtyard\textquotedblright{} (EtymDB), from Greek \emph{plateîa} \textquotedblleft{}broad\textquotedblright{}.
\end{cousinweb}

\begin{grammarlens}[title={What you've built}]
One more word for telling, explaining and describing at B1.
\end{grammarlens}

\subsection*{Your turn}
\begin{itemize}
\raggedright
  \item \emph{Say it:} \emph{la platea}
  \item \emph{Say it:} \emph{la platea}, once more
  \item \emph{Say it:} a sentence of your own with \emph{la platea}
  \item \emph{Recall:} say the Italian for stage, then the Italian for curtain, then say \emph{la platea} again
\end{itemize}

\begin{tcolorbox}[breakable,colback=teal!4,colframe=teal!35!black,title={Before you move on}]
What does la platea mean? (\textquotedblleft{}Stalls; audience\textquotedblright{}.) Say it once more.
\end{tcolorbox}
\section[il loggione]{il loggione — gallery, the gods (top seats)}

\label{lesson:IT-C519-loggione}



\begin{quote}
Before the new one: say the Italian for curtain, then the Italian for stalls.
\end{quote}

\subsection*{You'll want to know: il loggione}
\textbf{il loggione} — \textquotedblleft{}gallery, the gods (top seats)\textquotedblright{}.

\begin{quote}
Dal loggione si vede poco, ma costa meno.
\end{quote}

\begin{quote}
\emph{From the gods you see little, but it costs less.}
\end{quote}

The cheap seats at the top; the people there are \emph{i loggionisti}.

\begin{cousinweb}
\emph{loggione} is a big \emph{loggia} \textquotedblleft{}a gallery\textquotedblright{}, which came from Old French \emph{loge}, from Frankish \emph{laubija} \textquotedblleft{}a pergola\textquotedblright{} (EtymDB).
\end{cousinweb}

\begin{grammarlens}[title={What you've built}]
One more word for telling, explaining and describing at B1.
\end{grammarlens}

\subsection*{Your turn}
\begin{itemize}
\raggedright
  \item \emph{Say it:} \emph{il loggione}
  \item \emph{Say it:} \emph{il loggione}, once more
  \item \emph{Say it:} a sentence of your own with \emph{il loggione}
  \item \emph{Recall:} say the Italian for curtain, then the Italian for stalls, then say \emph{il loggione} again
\end{itemize}

\begin{tcolorbox}[breakable,colback=teal!4,colframe=teal!35!black,title={Before you move on}]
What does il loggione mean? (\textquotedblleft{}Gallery, the gods (top seats)\textquotedblright{}.) Say it once more.
\end{tcolorbox}
\section[l'applauso]{l'applauso — applause, round of applause}

\label{lesson:IT-C519-applauso}



\begin{quote}
Before the new one: say the Italian for stalls, then the Italian for gallery.
\end{quote}

\subsection*{You'll want to know: l'applauso}
\textbf{l'applauso} — \textquotedblleft{}applause, round of applause\textquotedblright{}.

\begin{quote}
Alla fine c'è stato un lungo applauso.
\end{quote}

\begin{quote}
\emph{At the end there was long applause.}
\end{quote}

\emph{Fare un applauso} — to give a round of applause.

\begin{cousinweb}
\emph{applauso} is a learned borrowing from Latin \emph{applausus} \textquotedblleft{}a clapping\textquotedblright{} (EtymDB), from \emph{applaudere} \textquotedblleft{}to strike upon, to clap\textquotedblright{}.
\end{cousinweb}

\begin{grammarlens}[title={What you've built}]
One more word for telling, explaining and describing at B1.
\end{grammarlens}

\subsection*{Your turn}
\begin{itemize}
\raggedright
  \item \emph{Say it:} \emph{l'applauso}
  \item \emph{Say it:} \emph{l'applauso}, once more
  \item \emph{Say it:} a sentence of your own with \emph{l'applauso}
  \item \emph{Recall:} say the Italian for stalls, then the Italian for gallery, then say \emph{l'applauso} again
\end{itemize}

\begin{tcolorbox}[breakable,colback=teal!4,colframe=teal!35!black,title={Before you move on}]
What does l'applauso mean? (\textquotedblleft{}Applause, round of applause\textquotedblright{}.) Say it once more.
\end{tcolorbox}
